% TRASER-F: прогнозирование графов сцены на PSG4D-HOI.
% Сборка: tectonic traser_f_report.tex
\documentclass[11pt,a4paper]{article}

\usepackage{fontspec}
\setmainfont{cmunrm.otf}[BoldFont=cmunbx.otf,ItalicFont=cmunti.otf,BoldItalicFont=cmunbi.otf]
\setsansfont{cmunss.otf}[BoldFont=cmunsx.otf]
\setmonofont{cmuntt.otf}[BoldFont=cmuntb.otf]
\usepackage[english,russian]{babel}
\usepackage[margin=2.2cm]{geometry}
\usepackage{amsmath,amssymb}
\usepackage{booktabs,tabularx,array,multirow}
\usepackage{xcolor}
\usepackage{enumitem}
\usepackage{listings}
\usepackage{tikz}
\usetikzlibrary{arrows.meta,positioning,fit,calc}
\usepackage{xurl}
\usepackage[hidelinks]{hyperref}
\emergencystretch=3em

\selectlanguage{russian}
\setlist{itemsep=2pt,topsep=4pt}
\definecolor{codebg}{RGB}{246,246,246}
\definecolor{accent}{RGB}{30,90,160}
\lstset{basicstyle=\ttfamily\small,backgroundcolor=\color{codebg},frame=single,
  rulecolor=\color{black!15},breaklines=true,columns=fullflexible,keepspaces=true,
  extendedchars=true,xleftmargin=4pt,xrightmargin=4pt}
\newcommand{\code}[1]{\texttt{#1}}
\newcolumntype{L}{>{\raggedright\arraybackslash}X}

\title{\textbf{TRASER-F: прогнозирование графов сцены}\\[2pt]
  \large Тройки и локализация объектов на PSG4D-HOI}
\author{Репозиторий \href{https://github.com/lve-gh/SVG2/tree/forecasting}{\code{lve-gh/SVG2}}, ветка \code{forecasting}}
\date{18 октября 2026}

\begin{document}
\maketitle

\begin{abstract}
TRASER (ECCV 2026) строит граф сцены по уже увиденному видео. Мы добавили к нему прогноз:
по первым $H$ секундам клипа и маскам объектов модель предсказывает отношения
(субъект, предикат, объект) на следующие $F$ секунд и положение каждого объекта, $H,F\in[1,5]$~с.
Документ описывает выбранную архитектуру TRASER-F и почему выбрана именно она, обзор
современных подходов, протокол оценки, написанный код, что уже проверено и что остаётся
сделать на сервере. Обучение на полных весах ещё не запускалось.
\end{abstract}

\tableofcontents

%======================================================================
\section{Задача и данные}

\paragraph{Постановка.}
\begin{center}
\begin{tabularx}{\linewidth}{@{}lL@{}}
\toprule
Вход & кадры интервала $[T-H,\,T)$ с частотой 2~кадра/с, маски объектов на этих кадрах, классы объектов \\
Выход & \code{\{"relationships": [[s, pred, o, [[t0, t1]]], ...]\}} для секунд $H,\dots,H+F-1$ клипа;
        для каждого объекта и каждой будущей секунды~--- бокс $(c_x,c_y,w,h)$ и видимость \\
Объекты & видимые в наблюдаемой части; новые объекты в будущем не появляются (как в SGA) \\
\bottomrule
\end{tabularx}
\end{center}
Это соответствует режимам GAGS/PGAGS из scene graph anticipation (объекты даны, прогнозируются
отношения) плюс локализация из scene graph \emph{forecasting} (FDSG).

\paragraph{PSG4D-HOI.} Часть PSG4D, построенная на HOI4D (без GTA-части):
\begin{itemize}
  \item 2973 эгоцентричных видео по 20~с, 15~кадров/с (300 кадров), $1920\times1080$;
  \item 46 классов объектов и около 15 предикатов: \emph{reachout, grasp, pickup, putdown,
        carry, carrywithbothhands, open, close, push, pull, press, cut, dump, binding};
  \item субъект почти всегда \code{hand}; все исполнители (H1--H4)~--- люди, роборук нет.
\end{itemize}
В HOI4D действия размечены отрезками времени (\code{action/color.json}), маски объектов~---
цветные PNG (\code{2Dseg/shift\_mask}). Из этого же PSG4D получал свои отношения
(«hand \textless действие\textgreater{} \textless объект\textgreater»), поэтому
конвертер умеет собирать разметку и напрямую из HOI4D (раздел~\ref{sec:code}).

\paragraph{Особенность данных.} Действия короткие: grasp $\approx 0{,}45$~с, pickup
$\approx 0{,}4$~с. В одну секунду часто попадают 2--3 предиката подряд, и их порядок внутри
секунды теряется. Это учтено в открытых вопросах (раздел~\ref{sec:open}).

%======================================================================
\section{Обзор современных подходов}

\begin{center}
\small
\begin{tabularx}{\linewidth}{@{}p{2.6cm}p{1.6cm}Lp{0.9cm}L@{}}
\toprule
Метод & Где & Как прогнозирует & Боксы & Что взяли в TRASER-F \\
\midrule
STTran+/++, DSG-DETR+/++ & базовые в SceneSayer & авторегрессионный трансформер & нет & авторегрессия накапливает ошибку \\
SceneSayer ODE/SDE & ECCV'24 & динамика признаков пар объектов в латентном пространстве (NeuralODE/SDE) & да & \textbf{прогноз в латентах}, декодирование теми же головами \\
ImparTail & CVPR'25 & SceneSayer/STTran + curriculum с маскированием неполной разметки & --- & устойчивость к пропускам в разметке \\
FDSG & 2025 & NeuralSDE на запросах DETR, бокс = тождество + поправка & \textbf{да} & \textbf{бокс = последний + $\Delta$}; метрики SGF с IoU $\ge 0{,}5$ \\
HyperGLM & CVPR'25 & гиперграф + мультимодальная LLM & нет & LLM-декодер конкурентен в SGA \\
OOTSM & 2025 & LLM по текстовым графам: сначала объекты, потом отношения & нет & LLM даёт больше всего в mR (+21,9 mR@50) \\
FORESCENE & 2025 & латентная диффузия графов & нет & вероятностный прогноз~--- следующий шаг \\
\bottomrule
\end{tabularx}
\end{center}

Ни один из методов не работает с токенами траекторий объектов внутри видео-LLM.
TRASER-F переносит идею SceneSayer/FDSG (латентная динамика, условие на горизонт, поправка
к последнему состоянию) в собственное представление TRASER и оставляет LLM TRASER
декодером, как в HyperGLM/OOTSM.

%======================================================================
\section{Архитектура TRASER-F}

\subsection{Исходный TRASER}
Qwen2.5-VL-3B и два Perceiver-ресемплера. Токены ViT отбираются по покрытию маской объекта;
для каждого объекта строится блок: OTR (вся траектория, 32 латента) и TWR (32 латента на
каждое временное окно). LLM читает блоки объектов и выдаёт граф сцены в JSON.
Мы берём частоту 2~кадра/с: при временном патче~2 один грид ViT равен одной секунде,
т.\,е. одному окну TWR.

\subsection{Схема}

\begin{center}
\begin{tikzpicture}[
  box/.style={draw, rounded corners=2pt, align=center, font=\small, minimum height=8mm, inner sep=4pt},
  new/.style={box, fill=accent!12, draw=accent, thick},
  old/.style={box, fill=black!4},
  arr/.style={-{Stealth[length=2.2mm]}, thick}]
  \node[old] (v) {кадры истории\\$[T-H,T)$};
  \node[old, right=6mm of v] (vit) {ViT\\Qwen2.5-VL};
  \node[old, right=6mm of vit] (sel) {отбор токенов\\по маскам};
  \node[old, right=6mm of sel] (res) {OTR + TWR\\на объект};
  \coordinate (mid) at ($(v.south west)!0.5!(res.south east)$);
  \node[new, below=11mm of mid, anchor=north, text width=9.5cm] (fc) {\textbf{Прогнозатор латентов} (новый, $\approx$24M)\\
     токены: (объект, окно истории) + OTR + запрос (объект, будущее окно)\\
     энкодер по всем объектам $\to$ декодер: 32 латента на запрос\\
     латенты $=$ последние $+\,\Delta$, \quad бокс $=$ последний $+\,\Delta$, видимость};
  \node[old, below=10mm of fc, text width=11.5cm, font=\ttfamily\footnotesize] (blk) {%
     <obj\_traj\_start> Object k (bottle): <vision\_start> [OTR]\\
     <0 - 1 sec>[TWR] ... <H-1 - H sec>[TWR]\\
     \textcolor{accent}{<forecast H - H+1 sec>[прогноз] ... <forecast H+F-1 - H+F sec>[прогноз]}\\
     <vision\_end> <obj\_traj\_end> \quad $\times$ каждый объект};
  \node[old, below=8mm of blk] (llm) {LLM TRASER $\to$ будущие отношения в JSON};
  \draw[arr] (v) -- (vit); \draw[arr] (vit) -- (sel); \draw[arr] (sel) -- (res);
  \draw[arr] (res.south) |- ($(fc.north)+(0,5mm)$) -- (fc.north);
  \draw[arr] (fc) -- (blk); \draw[arr] (blk) -- (llm);
  \draw[arr, dashed] (res.east) -- ++(5mm,0) |- node[pos=0.25, right, font=\footnotesize, align=left] {латенты\\истории} (blk.east);
\end{tikzpicture}
\end{center}

\subsection{Прогнозатор латентов}
Класс \code{ObjectLatentForecaster} (\code{forecast/dynamics.py}).
\begin{itemize}
  \item \textbf{Токены.} Для каждого объекта: по токену на окно истории
        (усреднённые латенты TWR после LayerNorm и проекции + бокс окна), токен OTR и по запросу на
        каждое будущее окно. К каждому токену добавляются эмбеддинги абсолютного времени, типа
        токена (история / запрос / OTR) и индекса объекта.
  \item \textbf{Энкодер.} Трансформер (4 слоя, $d=512$, 8 голов, pre-norm) по \emph{всем}
        токенам клипа: объекты видят друг друга.
  \item \textbf{Декодер.} 2 слоя: 32 обучаемых запроса, сдвинутых выходом энкодера для
        данного (объект, окно), смотрят на последние наблюдённые латенты объекта и дают
        $R=32$ латента будущего окна.
  \item \textbf{Выходы.}
        $\hat z_{t} = z_{\text{last}} + W_{\text{out}}\,h_t$,\quad
        $\hat b_t = \operatorname{clip}(b_{\text{last}} + \Delta b_t)$,\quad логит видимости.
        $W_{\text{out}}$ и голова $\Delta b$ инициализированы нулями: до обучения модель
        \emph{в точности} равна базовому copy-last.
  \item Около 23,7M параметров при $d_{\text{LLM}}=2048$.
\end{itemize}

\subsection{Почему эта архитектура}
\begin{enumerate}
  \item \textbf{Прогноз в латентах TRASER.} Будущие слоты заполняются латентами того же вида,
        что LLM уже умеет читать, поэтому декодер графа TRASER переиспользуется как есть.
  \item \textbf{Сразу весь горизонт, без авторегрессии.} Каждое будущее окно предсказывается
        напрямую по истории (запрос = объект + абсолютное время), ошибка на 5~с не накапливается.
  \item \textbf{Взаимодействие объектов.} Рука, тянущаяся к бутылке, меняет будущее бутылки.
  \item \textbf{Старт с copy-last.} Обучение отходит от инерции только там, где это помогает;
        модель не может стать хуже базового метода из-за плохой инициализации.
  \item \textbf{Бесплатная плотная разметка.} На обучении будущие кадры известны: латенты,
        которые ресемплер строит по ним, служат целями дистилляции без дополнительной разметки.
  \item \textbf{Нет утечки будущего.} ViT Qwen2.5-VL работает внутри одного временного грида
        (\code{cu\_seqlens}), поэтому кодирование всего клипа даёт те же признаки истории, что
        кодирование одной истории (проверено тестом). При оценке модели подаются только
        наблюдаемые кадры.
\end{enumerate}

\subsection{Функция потерь и обучение}
\[
\begin{aligned}
\mathcal L = {} & \mathcal L_{\text{LM}}(\text{будущий JSON})
  + \lambda_{\text{lat}}\,\bigl\lVert \mathrm{LN}(\hat z) - \mathrm{LN}(\operatorname{sg}(z^{\text{teacher}}))\bigr\rVert_2^2 \\
  & + \lambda_{\text{box}}\,\bigl(5\,\mathcal L_{1} + 2\,\mathcal L_{\text{GIoU}} + \mathcal L_{\text{BCE}}^{\text{vis}}\bigr),
\qquad \lambda_{\text{lat}}=\lambda_{\text{box}}=1.
\end{aligned}
\]
Цели $z^{\text{teacher}}$~--- латенты TWR реальных будущих окон, посчитанные без градиента.

\begin{center}
\begin{tabularx}{\linewidth}{@{}lp{3.6cm}p{2.6cm}L@{}}
\toprule
Этап & Обучается & LR & Зачем \\
\midrule
A, 20 эпох & только прогнозатор & $10^{-4}$ & выучить динамику против замороженного TRASER: дёшево, учитель стабилен \\
B, 10 эпох & LLM \newline merger \newline ресемплеры \newline прогнозатор
           & $10^{-5}$ \newline $2\cdot10^{-5}$ \newline $5\cdot10^{-5}$ \newline $5\cdot10^{-5}$ & адаптировать декодер к прогнозу \\
\bottomrule
\end{tabularx}
\end{center}
DeepSpeed ZeRO-3 с offload, bf16, накопление градиента 4, косинусный LR.
Клипы собираются на лету: случайные $H,F\in[1,5]$; в 70\,\% случаев якорь выбирается так,
чтобы в будущем окне начиналось отношение; клипы с пустым будущим сохраняются в 10\,\% случаев.

\paragraph{Режимы} (\code{--forecast\_mode}, один флаг): \code{latent}~--- основная модель;
\code{copy}~--- в будущие слоты копируются последние латенты; \code{oracle}~--- настоящие
будущие латенты (потолок); \code{prompt}~--- будущих слотов нет, LLM прогнозирует только по истории.

%======================================================================
\section{Протокол оценки}

\paragraph{Данные.} PSG4D-HOI, разбиение \emph{по экземплярам объектов} (\code{C*/N*} в пути
HOI4D), 75/10/15: тестовые объекты в обучении не встречаются.

\paragraph{Клипы.} Каждую секунду каждого тестового видео, для $(H,F)\in\{1,3,5\}^2$.
Основной результат~--- $H=3$, $F=3$. Все методы, включая базовые без обучения, оцениваются
на одних и тех же клипах (общая функция \code{clip\_targets}).

\paragraph{Метрики} (\code{forecast/metrics.py}), по будущим секундам:
\begin{center}
\begin{tabularx}{\linewidth}{@{}lL@{}}
\toprule
R, mR & полнота четвёрок (s, pred, o, секунда); mR~--- среднее по предикатам \\
P, F1 & точность, чтобы не окупалось перечисление всех предикатов \\
R@$k$ & полнота на $k$-й секунде вперёд (как быстро падает качество) \\
mIoU@$k$, ADE, FDE & прогноз бокса против реального \\
R\textsubscript{loc} & как в SGF: тройка верна \emph{и} боксы обоих объектов с IoU $\ge 0{,}5$ \\
\bottomrule
\end{tabularx}
\end{center}

\paragraph{Базовые методы} (\code{forecast/baselines.py}). Для троек: copy-last и марковская
цепочка переходов предикатов; для боксов: copy и постоянная скорость. Им даётся истинный
наблюдаемый граф, т.\,е. больше, чем видит модель, поэтому обойти их~--- содержательный рубеж.

\paragraph{ImparTail / SceneSayer.} Их код привязан к Action Genome (боксы по кадрам,
признаки Faster R-CNN). Два сравнения:
\begin{enumerate}
  \item \emph{На PSG4D-HOI:} обучить SceneSayerSDE (с curriculum ImparTail и без) на нашем
        разбиении в режиме PGAGS/GAGS, боксы из наших масок, 1 кадр/с; оценка
        \code{recall\_at\_k} на наших четвёрках.
  \item \emph{На Action Genome:} прогнать TRASER-F с прямоугольными масками из разметки AG по
        протоколу SGA ($F\in\{0{,}3;0{,}5;0{,}7;0{,}9\}$, R@10/20/50 With/No constraint) и
        сравнить с опубликованными SceneSayerSDE, ImparTail, FDSG.
\end{enumerate}

\paragraph{Гипотезы и абляции.} Всё с одинаковыми данными, расписанием и оценкой; меняется один флаг.
\begin{center}
\begin{tabularx}{\linewidth}{@{}lLl@{}}
\toprule
ID & Гипотеза & Запуск \\
\midrule
B0 & инерция~--- сильный базовый метод & \code{baselines.py}, \code{--mode copy} \\
H1 & LLM может прогнозировать по одной истории & \code{MODE=prompt} \\
H2 & латентная динамика лучше инерции (основная модель) & \code{MODE=latent} \\
H2a & работает именно дистилляция латентов & \code{LAMBDA\_LATENT=0} \\
H2b & боксы помогают отношениям & \code{LAMBDA\_BOX=0} \\
UB & потолок при идеальной динамике & \code{evaluate.py --mode oracle} \\
H3 & длина истории важна & одна модель, $H=1\dots5$ \\
\bottomrule
\end{tabularx}
\end{center}
Сравнение \code{oracle}/\code{latent}/\code{copy} разделяет ошибку на «динамику» и
«декодирование»: если \code{oracle} ненамного лучше \code{copy}, узкое место в декодере.

%======================================================================
\section{Что написано}\label{sec:code}

Весь код~--- новые файлы; исходный код TRASER не изменён.
\begin{center}
\small
\begin{tabularx}{\linewidth}{@{}p{6.7cm}L@{}}
\toprule
Файл & Содержание \\
\midrule
\multicolumn{2}{@{}l}{\code{traser/traser\_train/forecast/}} \\
\quad\code{data\_forecast.py} & формат разметки, нарезка клипов, боксы из RLE-масок, \code{ForecastDataset}, коллатор \\
\quad\code{dynamics.py} & прогнозатор латентов, лоссы дистилляции и боксов (L1, GIoU, BCE) \\
\quad\code{arrangement.py} & сборка входа LLM: блоки объектов с будущими слотами, позиции, \code{rope\_deltas} \\
\quad\code{step.py} & отбор токенов по маскам для истории и будущего, метки времени \\
\quad\code{modeling\_forecast.py} & \code{TRASERForecast(TRASER)}, загрузка из чекпойнта TRASER, инициализация под ZeRO-3 \\
\quad\code{trainer\_forecast.py} & \code{ForecastTrainer}: группы LR, суммарный лосс, логирование составляющих \\
\quad\code{train\_forecast.py} & точка входа обучения, аргументы \\
\quad\code{metrics.py} & разбор JSON, R/P/F1/mR/R@$k$, mIoU, ADE/FDE, R\textsubscript{loc} \\
\quad\code{baselines.py} & copy-last, марковская цепочка, боксы copy/velocity \\
\quad\code{evaluate.py} & оценка: только кадры истории, генерация, метрики, jsonl \\
\midrule
\code{traser/data/prepare\_psg4d\_hoi.py} & конвертер: \code{psg4d} (JSON PSG4D-HOI) и \code{hoi4d} (действия HOI4D) \\
\code{traser/scripts/train\_forecast.sh} & этапы A/B через DeepSpeed \\
\code{traser/tests/test\_forecast\_smoke.py} & smoke-тест на крошечной случайной модели (CPU) \\
\code{docs/FORECAST.md} & описание на английском \\
\code{docs/report/} & этот отчёт (\code{.tex} и \code{.pdf}) \\
\bottomrule
\end{tabularx}
\end{center}

\paragraph{Формат разметки.}
\begin{lstlisting}
{"video": ".../image.mp4", "mask_json": ".../<seq>_rle.json", "fps": 15,
 "num_frames": 300, "objects": [{"id": 0, "name": "toy car"}, {"id": 1, "name": "right hand"}],
 "relationships": [[1, "grasp", 0, [[45, 52], [180, 190]]], ...]}
\end{lstlisting}
\code{mask\_json}~--- список по кадрам, в каждом COCO-RLE для каждого объекта; интервалы
отношений~--- в исходных кадрах.

%======================================================================
\section{Что проверено}

\paragraph{Smoke-тест} (\code{python traser/tests/test\_forecast\_smoke.py}; случайная
Qwen2.5-VL с $d=64$, синтетический клип 8~с: рука движется к бутылке)~--- проходит целиком:
\begin{enumerate}[label=\arabic*.]
  \item датасет собирает клипы с нужным числом окон истории и будущего;
  \item утечки будущего через ViT нет;
  \item при инициализации прогноз в точности равен copy-last, в том числе
        (3b) при загрузке чекпойнта TRASER без прогнозатора;
  \item градиенты доходят до прогнозатора; два шага \code{ForecastTrainer.train()};
  \item оценка: генерация, разбор, метрики, базовые методы.
\end{enumerate}
Найден и исправлен баг: стандартная инициализация HF не трогает \code{in\_proj\_weight} у
\code{nn.MultiheadAttention}, и при загрузке из чекпойнта TRASER эти веса остались бы
неинициализированными.

\paragraph{Реальное видео HOI4D} (категория C1, игрушечная машинка; объекты toy car и right
hand, 12 отрезков отношений, разметка через конвертер \code{hoi4d}): сборка клипа 3,7~с,
грид $6\times40\times72$, длина собранной последовательности 341 токен, три шага обучения
на крошечной модели прошли.

\begin{center}
\begin{tabular}{@{}llrrr@{}}
\toprule
$H/F$ & Метод & R, \% & mR, \% & P, \% \\
\midrule
3/3 & copy-last & 34,4 & 25,2 & --- \\
3/3 & Markov    & 46,9 & --- & 66,7 \\
5/5 & copy-last & 20,2 & --- & --- \\
5/5 & Markov    & 34,5 & --- & --- \\
\bottomrule
\end{tabular}
\end{center}
\textbf{Это не результат}, а проверка, что конвейер работает: видео одно, а марковская цепочка
обучена на том же видео. Для боксов постоянная скорость хуже copy в эгоцентричном видео
(mIoU@1с 46 против 60): камера движется вместе с головой.

\paragraph{Не проверено:} запуск с реальными весами 3B, DeepSpeed ZeRO-3 и flash-attn;
конвертер \code{psg4d} на настоящих файлах PSG4D-HOI (формат восстановлен по скрипту
визуализации авторов, файлы на SharePoint требуют входа).

%======================================================================
\section{Запуск}

\begin{lstlisting}
git clone -b forecasting https://github.com/lve-gh/SVG2 && cd SVG2

# 1. Данные: PSG4D-HOI или запасной вариант из действий HOI4D
python traser/data/prepare_psg4d_hoi.py psg4d --psg4d_json HOI/psg4d_hoi.json \
    --categories HOI/categories.json --hoi4d_root HOI/HOI4D_release \
    --mask_root HOI/HOI4D_mask --out_dir traser/data
python traser/data/prepare_psg4d_hoi.py hoi4d --hoi4d_root HOI4D_release \
    --annotation_root HOI4D_annotations --instructions HOI4D-Instructions \
    --out_dir traser/data

# 2. Проверка окружения и базовые методы
python traser/tests/test_forecast_smoke.py
cd traser && python -m traser_train.forecast.baselines \
    --train data/psg4d_hoi_train.json --test data/psg4d_hoi_test.json \
    --history_sec 3 --horizon_sec 3

# 3. Обучение (NUM_GPUS, MODE, LAMBDA_LATENT, LAMBDA_BOX, EXTRA_ARGS)
STAGE=A bash traser/scripts/train_forecast.sh
STAGE=B bash traser/scripts/train_forecast.sh

# 4. Оценка
python -m traser_train.forecast.evaluate --model checkpoints/traser_f_latent_stageB \
    --annotation data/psg4d_hoi_test.json --history_sec 3 --horizon_sec 3 \
    --mode latent --output preds.jsonl
\end{lstlisting}

%======================================================================
\section{Открытые вопросы}\label{sec:open}

\begin{itemize}
  \item \textbf{Файлы PSG4D-HOI.} Нужны \code{psg4d\_hoi.json} и \code{categories.json}
        (SharePoint NTU, нужен вход). До этого~--- разметка из действий HOI4D.
  \item \textbf{3D.} Боксы сейчас 2D (\code{forecaster\_box\_dim=4}). Для 3D-версии TRASER
        достаточно задать размерность 3D-бокса и считать цели по глубине; остальной
        прогнозатор от этого не зависит. Нужен код командной 3D-адаптации.
  \item \textbf{Временное разрешение.} Если метрики по секундам упрутся в потолок~--- слоты
        по 0,5~с (\code{sample\_fps 4}, полусекундные метки); прогнозатор не меняется.
  \item \textbf{Многовариантное будущее.} Прогнозатор детерминированный. Если H2 работает, но
        R@$k$ быстро падает с $k$,~--- стохастический латент (SDE, как в SceneSayer/FDSG) и
        оценка best-of-N.
  \item \textbf{Сервер.} Нужны число и тип GPU (\code{nvidia-smi}), чтобы подобрать
        \code{NUM\_GPUS}, батч и конфигурацию DeepSpeed.
\end{itemize}

%======================================================================
\section*{Источники}
\begin{itemize}
  \item TRASER / SVG2: \url{https://github.com/uwGZQ/SVG2}
  \item PSG4D: \url{https://arxiv.org/abs/2405.10305}, \url{https://github.com/Jingkang50/PSG4D}
  \item SceneSayer (Towards Scene Graph Anticipation): \url{https://arxiv.org/abs/2403.04899}
  \item ImparTail: \url{https://github.com/rohithpeddi/ImparTail}
  \item FDSG: \url{https://arxiv.org/abs/2506.01487}
  \item OOTSM: \url{https://arxiv.org/abs/2509.05661}
  \item HyperGLM (CVPR 2025): \url{https://openaccess.thecvf.com/content/CVPR2025/papers/Nguyen_HyperGLM_HyperGraph_for_Video_Scene_Graph_Generation_and_Anticipation_CVPR_2025_paper.pdf}
  \item FORESCENE: \url{https://arxiv.org/abs/2503.06182}
\end{itemize}

\end{document}
